\title{Direct Preference Optimization: Your Language Model is Secretly a Reward Model}

\begin{abstract}
Although large-scale unsupervised language models (LMs) acquire extensive world knowledge and some reasoning abilities, precisely controlling their behavior remains challenging due to the entirely unsupervised nature of their training. Current approaches to achieving such steerability involve collecting human judgments about the relative quality of model outputs and fine-tuning the unsupervised LM to align with these preferences, often using reinforcement learning from human feedback (RLHF). Nevertheless, RLHF is a complicated and frequently unstable process, which first fits a reward model that captures human preferences, and then fine-tunes the large unsupervised LM via reinforcement learning to optimize this estimated reward while avoiding significant divergence from the original model. In this work, we present a novel parameterization of the reward model in RLHF that allows for direct extraction of the optimal policy in closed form, enabling us to address the standard RLHF problem using only a straightforward classification loss. The resulting method, termed \textit{Direct Preference Optimization} (DPO), is stable, effective, and computationally efficient, removing the necessity for sampling from the LM during fine-tuning or extensive hyperparameter tuning. Our experimental results demonstrate that DPO can fine-tune LMs to better align with human preferences compared to existing techniques. Importantly, fine-tuning with DPO outperforms PPO-based RLHF in controlling the sentiment of generated outputs, and matches or improves quality in summarization and single-turn dialogue tasks, while being considerably simpler to implement and train.
\end{abstract}

\section{Introduction}  
Large-scale unsupervised language models (LMs) trained on vast datasets develop remarkable capabilities~\citep{chowdhery2022palm, brown2020language, touvron2023llama,bubeck2023sparks}. Nonetheless, these models are trained on human-generated data that reflects a wide range of objectives, priorities, and skill levels. Some of these objectives and skills may not be desirable to replicate; for instance, while we want an AI coding assistant to \textit{recognize} common programming errors for correction purposes, when generating code, we prefer to steer the model toward the (often rare) instances of high-quality coding skills found in the training data. Likewise, we might want a language model to be \textit{aware} of a misconception commonly held by 50\% of people, but we certainly do not want it to assert this misconception as true in half of the responses to related queries. Put differently, it is essential to select the model’s \emph{preferred responses and behaviors} from its broad \textit{knowledge and capabilities} to create AI systems that are safe, effective, and controllable \citep{ouyang2022training}. Although current approaches commonly guide LMs to align with human preferences via reinforcement learning (RL), we demonstrate that the RL-based objective used in these methods can be precisely optimized using a straightforward binary cross-entropy loss, thereby significantly simplifying the preference learning process.

Broadly speaking, existing approaches instill desired behaviors into language models by using curated datasets of human preferences, which capture the types of behaviors humans deem safe and useful. This preference learning phase follows an initial large-scale unsupervised pre-training step on an extensive text corpus. While the simplest method for preference learning involves supervised fine-tuning on human demonstrations of high-quality outputs, the most effective techniques rely on reinforcement learning from human (or AI) feedback (RLHF/RLAIF; \citep{christiano2017deep,bai2022constitutional}). RLHF methods train a reward model on a dataset of human preferences and then employ RL to optimize the LM policy to generate responses that receive high reward, without deviating excessively from the original model. Although RLHF yields models with impressive conversational and coding performance, its pipeline is considerably more complex than supervised learning, requiring training multiple LMs and sampling from the LM policy during training, which entails substantial computational expense.

In this paper, we demonstrate a method to directly optimize a language model to align with human preferences, without relying on explicit reward modeling or reinforcement learning. We introduce \textit{Direct Preference Optimization} (DPO), an algorithm that implicitly optimizes the same objective as current RLHF algorithms (reward maximization constrained by KL-divergence) but is simpler to implement and more straightforward to train. Conceptually, the DPO update boosts the relative log probability of preferred responses over dispreferred ones, while incorporating a dynamic, per-example importance weighting that avoids the model collapse we observe when using a naive probability ratio objective. Similar to existing approaches, DPO depends on a theoretical preference model (e.g., the Bradley-Terry model; \cite{bradley1952rankanalysis}) to quantify how well a candidate reward function fits empirical preference data. However, unlike current methods which employ the preference model to formulate a preference loss for training a reward model and subsequently optimize a policy against that learned reward, DPO applies a change of variables to define the preference loss directly as a function of the policy. Given a dataset of human preferences over model outputs, DPO can thus optimize a policy via a simple binary cross entropy objective, yielding the optimal policy for an implicit reward function inferred from the preference data.

Our primary contribution is Direct Preference Optimization (DPO), a straightforward RL-free method for training language models from preferences. Our empirical results indicate that DPO performs at least as well as existing techniques, including PPO-based RLHF, in learning from preferences for tasks such as sentiment control, summarization, and dialogue generation, using language models with up to 6B parameters.

\section{Related Work}

Self-supervised language models with growing scale have demonstrated the ability to perform certain tasks in a zero-shot manner \citep{radford2019language} or through few-shot prompting \citep{gpt3,megatron,chowdhery2022palm}. Nevertheless, their effectiveness on downstream tasks and alignment with user goals can be substantially enhanced by fine-tuning on datasets containing instructions paired with human-generated completions \citep{mishra-etal-2022-cross,sanh2022multitask,chung2022scaling,thoppilan2022lamda}. This process, known as `instruction-tuning,' allows LLMs to generalize to instructions beyond those seen during tuning and generally improves their usability \citep{chung2022scaling}. Although instruction tuning has proven successful, collecting \textit{relative} human judgments on response quality is often simpler than acquiring expert demonstrations. Consequently, later studies have fine-tuned LLMs using datasets reflecting human preferences, thereby enhancing performance in areas such as translation \citep{kreutzer-etal-2018-reliability}, summarization \citep{stiennon2022learning,ziegler2020finetuning}, storytelling \citep{ziegler2020finetuning}, and instruction adherence \citep{ouyang2022training,ramamurthy2023is}. These approaches first train a neural network reward model to align with the preference dataset under models like the Bradley-Terry model \citep{bradley1952rankanalysis}, then fine-tune the language model to maximize this reward via reinforcement learning techniques, typically employing REINFORCE \citep{williams1992reinforce}, proximal policy optimization (PPO; \cite{schulman2017proximal}), or related variants \citep{ramamurthy2023is}. A closely related research direction exploits LLMs fine-tuned on human feedback for instruction following to generate synthetic preference data targeting specific qualities such as safety or harmlessness \citep{bai2022constitutional}, relying only on weak human supervision through textual rubrics guiding the LLM’s annotations. These approaches illustrate a merging of two research strands: one focused on training language models with reinforcement learning for diverse objectives \citep{Ranzato2015SequenceLT,paulus2018a,wu2018learning} and another on general frameworks for learning from human preferences \citep{christiano2017deep,kupcsik2018learning}. Despite the attractiveness of leveraging relative human preferences, applying reinforcement learning to fine-tune large language models remains a significant practical obstacle; this paper introduces a theoretically grounded method for optimizing relative preferences without using RL.

Outside the realm of language, policy learning from preferences has been explored within both bandit and reinforcement learning frameworks, leading to various proposed methods. Contextual bandit learning that utilizes preferences or rankings of actions instead of explicit rewards is referred to as a contextual dueling bandit (CDB; \cite{yue2012karmed,dudik2015contextual}). When absolute rewards are unavailable, theoretical treatments of CDBs replace the concept of an optimal policy with a \textit{von Neumann winner}, defined as a policy whose expected winning probability against \textit{any} other policy is at least 50\% \citep{dudik2015contextual}. Nonetheless, in the CDB paradigm, preference labels are provided online, whereas in human preference learning, the process typically relies on a fixed offline dataset of preference-labeled action pairs \citep{yan2022human}. Similarly, \textit{preference-based RL} (PbRL) learns from binary preferences produced by an \textit{unknown} latent `scoring' function rather than direct rewards \citep{BusaFekete2014,ruiz2023dueling}. Numerous PbRL algorithms exist, including those capable of leveraging off-policy preference data, but they generally entail first explicitly estimating the hidden scoring function (i.e., the reward model), followed by optimizing it \citep{jain2013learning,BusaFekete2014,christiano2017deep,sadigh2017active,kupcsik2018learning}. In contrast, we propose a single-stage policy learning method that directly optimizes the policy to align with the observed preferences.

\section{Preliminaries}\label{section:prelims}

We summarize the RLHF pipeline from \citeauthor{ziegler2020finetuning} and subsequent works \citep{stiennon2022learning, bai2022training, ouyang2022training}, which generally encompasses three stages: 1) supervised fine-tuning (SFT); 2) preference collection and reward learning; and 3) reinforcement learning optimization.

\textbf{SFT}: RLHF typically starts by fine-tuning a pre-trained language model on high-quality supervised data tailored to the target downstream tasks (e.g., dialogue, summarization), yielding a model denoted as $\pi^\text{SFT}$.

\textbf{Reward Modelling Phase}: In the second stage, the SFT model is prompted with inputs $x$ to generate pairs of responses $(y_1, y_2) \sim \pi^\text{SFT}(y \mid x)$. These pairs are then evaluated by human annotators, who indicate a preference for one response, written as $y_w \succ y_l \mid x$, where $y_w$ and $y_l$ represent the preferred and less preferred completions respectively among $(y_1, y_2)$. These preference labels are assumed to arise from an underlying latent reward function $r^*(y, x)$, which is not directly observable. Several models exist for capturing these preferences, with the Bradley-Terry (BT) model \cite{bradley1952rankanalysis} being a widely used approach (although the more general Plackett-Luce ranking models \citep{plackett1975analysis, luce2012individual} can also be integrated into this framework, especially if multiple ranked responses are available). The BT model expresses the human preference distribution $p^*$ as:

\begin{equation}\label{eq:bradley-terry}
    p^*(y_1\succ y_2 \mid x)=\frac{\exp\left(r^*(x, y_1)\right)}{\exp\left(r^*(x, y_1)\right) + \exp\left(r^*(x, y_2)\right)}.
\end{equation}

Given a fixed dataset of comparisons $\mathcal{D}=\bigl\{x^{(i)}, y_w^{(i)}, y_l^{(i)}\bigr\}_{i=1}^N$ drawn from $p^*$, one can define a reward model $r_{\phi}(x, y)$ and determine its parameters through maximum likelihood estimation. By formulating the problem as a binary classification task, the negative log-likelihood loss is expressed as:

\begin{equation}\label{eq:reward_model}
    \mathcal{L}_R(r_{\phi}, \mathcal{D}) = -\mathbb{E}_{(x, y_w, y_l)\sim \mathcal{D}}\bigl[\log \sigma(r_{\phi}(x, y_w)- r_{\phi}(x, y_l))\bigr]
\end{equation}

where $\sigma$ denotes the logistic function. For language models, the function $r_{\phi}(x, y)$ is typically initialized using the SFT model $\pi^\text{SFT}(y \mid x)$, augmented by an additional linear layer atop the final transformer layer that yields a scalar reward prediction \cite{ziegler2020finetuning}. To reduce variance in the reward function, previous work often normalizes the rewards such that $\mathbb{E}_{x,y\sim \mathcal{D}}\left[r_\phi(x, y)\right] = 0$ for every $x$.

\textbf{RL Fine-Tuning Phase}: During the reinforcement learning stage, the estimated reward function is employed to guide feedback to the language model. Following established methodologies~\citep{jaques2017sequence, jaques2020human}, the optimization problem is posed as

\begin{equation}\label{eq:RL}
\max_{\pi_{\theta}}  \mathbb{E}_{x\sim \mathcal{D}, y\sim \pi_{\theta}(y \mid x)}\bigl[r_{\phi}(x, y)\bigr] - \beta\mathbb{D}_{\textrm{KL}}\bigl[\pi_{\theta}(y\mid x)\mid \mid \pi_\text{ref}(y\mid x)\bigr],
\end{equation}

where $\beta$ regulates the allowed divergence from the baseline reference policy $\pi_\text{ref}$, usually the initial SFT model $\pi^\text{SFT}$. 
In practice, the language model policy $\pi_\theta$ is also initialized to $\pi^\text{SFT}$. This regularization is crucial as it restricts the policy from straying excessively from the distribution where the reward model remains valid, while preserving generative diversity and avoiding collapse to a narrow set of high-reward outputs. Because language generation is discrete, this objective is not differentiable and is commonly optimized via reinforcement learning algorithms. The prevailing method \citep{ziegler2020finetuning, stiennon2022learning, bai2022training, ouyang2022training} constructs a reward function defined as ${r(x, y) = r_{\phi}(x, y) -\beta (\log \pi_{\theta}(y\mid x) - \log \pi_\text{ref}(y\mid x))}$, which is then maximized using PPO \cite{schulman2017proximal}. 

\section{Direct Preference Optimization}\label{sec:DPO}

Motivated by the difficulties inherent in applying reinforcement learning techniques to large-scale tasks such as language model fine-tuning, we aim to develop a straightforward method for policy optimization that works directly with preference data. Contrary to prior RLHF approaches, which first learn a reward function and subsequently optimize it using RL, our method exploits a specific parameterization of the reward model that allows the optimal policy to be derived in closed form, eliminating the need for an RL training loop. As we will explain in detail next, the central insight lies in utilizing an analytical transformation from reward functions to their corresponding optimal policies, which enables reformulating a loss on reward functions into one defined over policies. This change-of-variables technique removes the requirement to fit an explicit, independent reward model while still optimizing according to established models of human preferences such as the Bradley-Terry model. Effectively, the policy network simultaneously embodies both the language model and

\textbf{Derivation of the DPO objective.} We begin with the same reinforcement learning objective as previous studies, Eq.~\ref{eq:RL}, defined under a general reward function $r$. Following earlier works~\citep{peters2007reinforcement, peng2019advantage, korbak2022reinforcement, go2023aligning}, it can be readily demonstrated that the optimal solution to the KL-constrained reward maximization problem in Eq.~\ref{eq:RL} is given by:

\begin{equation}\label{eq:op_policy}
    \pi_r(y\mid x) = \frac{1}{Z(x)}\pi_\text{ref}(y\mid x)\exp\left(\frac{1}{\beta}r(x, y)\right),
\end{equation}

where $Z(x) = \sum_{y}\pi_\text{ref}(y\mid x)\exp\left(\frac{1}{\beta}r(x, y)\right)$ represents the partition function. A full derivation is provided in Appendix \ref{app:derivation1}. Even if we substitute the maximum likelihood estimate $r_{\phi}$ for the true reward function $r^*$, estimating the partition function $Z(x)$ remains computationally expensive \citep{korbak2022reinforcement, go2023aligning}, limiting practical use of this formulation. Nonetheless, rearranging Eq.~\ref{eq:op_policy} allows us to express the reward function in terms of its corresponding optimal policy $\pi_r$, the reference policy $\pi_\text{ref}$, and the unknown partition function $Z(\cdot)$. Specifically, by taking the logarithm on both sides of Eq.~\ref{eq:op_policy} and performing some algebraic manipulation, we arrive at:

\begin{equation}\label{eq:main_eq}
    r(x,y) = \beta \log \frac{\pi_r(y\mid x)}{\pi_\text{ref}(y\mid x)} + \beta \log Z(x).
\end{equation}

This reparameterization can be applied to the true reward $r^*$ and its corresponding optimal policy $\pi^*$. Fortunately, the Bradley-Terry model depends solely on the difference in rewards between two outputs, that is, ${p^*(y_1 \succ y_2 \mid x) = \sigma(r^*(x, y_1) - r^*(x, y_2))}$. By substituting the reparameterized form from Eq.~\ref{eq:main_eq} for $r^*(x,y)$ into the preference model Eq.~\ref{eq:bradley-terry}, the partition function terms cancel out, enabling us to express the human preference probability purely in terms of the optimal policy $\pi^*$ and the reference policy $\pi_\text{ref}$. Thus, the optimal RLHF policy $\pi^*$ under the Bradley-Terry framework satisfies:

\begin{equation}\label{eq:objective}
    p^*(y_1\succ y_2 \mid x) = \frac{1}{1 + \exp\left(\beta \log \frac{\pi^*(y_2\mid x)}{\pi_\text{ref}(y_2\mid x)} - \beta \log \frac{\pi^*(y_1\mid x)}{\pi_\text{ref}(y_1\mid x)}\right)}.
\end{equation}

The full derivation is available in Appendix~\ref{app:derivation2}. Although Eq.~\ref{eq:objective} employs the Bradley-Terry model, similar derivations hold for the broader class of Plackett-Luce models~\citep{plackett1975analysis, luce2012individual}, as detailed in Appendix~\ref{app:plackett_luce_models}.

Having expressed the human preference likelihood in terms of the optimal policy rather than the reward function, we can now define a maximum likelihood objective for a parameterized policy $\pi_\theta$. Mirroring the reward modeling objective (Eq.~\ref{eq:reward_model}), the policy loss becomes:

\begin{equation}\label{eq:optimum_model}
    \mathcal{L}_\text{DPO}(\pi_{\theta}; \pi_\text{ref}) = -\mathbb{E}_{(x, y_w, y_l)\sim \mathcal{D}}\left[\log \sigma \left(\beta \log \frac{\pi_{\theta}(y_w\mid x)}{\pi_\text{ref}(y_w\mid x)} - \beta \log \frac{\pi_{\theta}(y_l

In this manner, we fit an implicit reward using a different parameterization, whose optimal policy is simply $\pi_\theta$. Additionally, since our approach corresponds to fitting a reparametrized Bradley-Terry model, it benefits from certain theoretical guarantees, such as consistency under appropriate assumptions regarding the preference data distribution \cite{bong2022generalized}. Section~\ref{sec:theory} further explores the theoretical aspects of DPO in comparison to related works.

\textbf{What effect does the DPO update have?} To gain a mechanistic insight into DPO, it is helpful to examine the gradient of the loss function $\mathcal{L}_\text{DPO}$. The gradient with respect to the parameters $\theta$ can be expressed as:

\begin{multline*}\label{eq:gradient}
    \nabla_\theta \mathcal{L}_\text{DPO}(\pi_\theta;\pi_\text{ref}) = \\ -\beta\mathbb{E}_{(x, y_w, y_l) \sim \mathcal{D}} \bigg[\underbrace{\sigma(\hat{r}_\theta(x, y_l) - \hat{r}_\theta (x, y_w))}_\text{larger weight when reward estimate is incorrect}\bigg[\underbrace{\nabla_\theta\log \pi(y_w \mid x)}_\text{boost likelihood of $y_w$} - \underbrace{\nabla_\theta\log\pi(y_l \mid x)}_\text{reduce likelihood of $y_l$}\bigg]\bigg],
\end{multline*}

where $\hat{r}_\theta(x, y) = \beta \log \frac{\pi_\theta(y \mid x)}{\pi_\text{ref}(y \mid x)}$ defines the reward implicitly through the language model $\pi_\theta$ and the reference model $\pi_\text{ref}$ (discussed further in Section~\ref{sec:theory}). Intuitively, the gradient of $\mathcal{L}_\text{DPO}$ increases the probability of preferred outputs $y_w$ while decreasing that of less preferred outputs $y_l$. Crucially, the samples are weighted based on the degree to which the implicit reward model $\hat{r}_\theta$ mistakenly ranks the less preferred outputs higher, scaled by $\beta$, reflecting the intensity of the KL regularization. Our empirical results highlight the importance of this weighting factor, as a naive implementation without it can lead to the language model collapsing (see Appendix Table~\ref{tab:unlikelihood_generations}).

\textbf{DPO overview.} The typical DPO process involves: 1) Sampling completions $y_1, y_2 \sim \pi_\text{ref}(\cdot \mid x)$ for each prompt $x$, then labeling these completions with human preferences to create the offline preference dataset $\mathcal{D} = \{x^{(i)}, y_w^{(i)}, y_l^{(i)}\}_{i=1}^N$; 2) optimizing the language model $\pi_\theta$ to minimize $\mathcal{L}_\text{DPO}$ given the reference model $\pi_\text{ref}$, dataset $\mathcal{D}$, and the target $\beta$. 

In practical scenarios, one often prefers to utilize publicly available preference datasets instead of generating new samples and collecting human preferences. Since these preference datasets are collected using $\pi^\text{SFT}$, we set $\pi_\text{ref} = \pi^\text{SFT}$ whenever it exists. If $\pi^\text{SFT}$ is unavailable, we initialize $\pi_\text{ref}$ by maximizing the likelihood of preferred completions ${(x, y_w)}$, specifically, ${\pi_\text{ref} = \argmax_{\pi}\mathbb{E}_{x, y_w \sim \mathcal{D}}\left[\log \pi(y_w \mid x)\right]}$. This strategy helps reduce the distribution shift between the unknown true reference distribution and the $\pi_\text{ref}$ employed by DPO. Additional information on implementation and hyperparameters is provided in Appendix~\ref{app:implementation}.

\section{Theoretical Analysis of DPO} In this section, we offer a deeper interpretation of the DPO method, supply theoretical justification, and connect the benefits of DPO to the challenges faced by actor-critic algorithms used in RLHF (such as PPO~\cite{schulman2017proximal}).

\label{sec:theory}

\subsection{Your Language Model Is Secretly a Reward Model} DPO manages to bypass the need for explicitly fitting a reward function and performing RL by optimizing a single maximum likelihood objective. Observe that the optimization objective Eq. \ref{eq:main_eq} is equivalent to a Bradley-Terry model where the reward is parameterized as $r^*(x, y) = \beta \log\frac{\pi^*_\theta(y \mid x)}{\pi_\text{ref}(y \mid x)}$. We optimize the parametric model $\pi_{\theta}$ in a way that corresponds directly to the reward model optimization in Eq. \ref{eq:reward_model} under a variable transformation. In this section, we develop the theoretical foundation for this reparameterization, demonstrate that it does not restrict the variety of reward models that can be learned, and show that it allows for exact recovery of the optimal policy. We begin by defining an equivalence relation between reward functions.

\begin{definition}
We define two reward functions $r(x, y)$ and $r'(x, y)$ as equivalent if and only if there exists a function $f$ such that ${r(x, y)-r'(x, y) = f(x)}$.     
\end{definition}

It is straightforward to verify that this relation is indeed an equivalence relation, which divides the set of reward functions into equivalence classes. We can then state the following two lemmas:

\begin{lemma}\label{lemma:same_prefrence} Within the Plackett-Luce framework, and specifically the Bradley-Terry model, two reward functions belonging to the same equivalence class produce the same preference distribution.
\end{lemma}

\begin{lemma}\label{lemma:same_policy}
    Reward functions in the same equivalence class generate the identical optimal policy for the constrained RL problem.
\end{lemma}

The proofs are simple and provided in Appendix \ref{app:lemma1}. The first lemma highlights a known under-specification problem inherent to the Plackett-Luce family of models \cite{plackett1975analysis}. Due to this ambiguity, additional identifiability constraints are typically necessary to secure guarantees on the MLE estimates from Eq. \ref{eq:reward_model} \cite{bong2022generalized}. The second lemma indicates that all reward functions within the same equivalence class yield the same optimal policy, so our ultimate goal is to recover any representative reward function from the optimal class. We establish the following theorem in Appendix~\ref{app:thm1}:

\begin{theorem}\label{thm:main}
    Given mild assumptions, every reward class consistent with the Plackett-Luce (and especially Bradley-Terry) models can be expressed by the reparameterization ${r(x, y) = \beta \log \frac{\pi(y\mid x)}{\pi_\text{ref}(y\mid x)}}$, for some model $\pi(y\mid x)$ and a specified reference model $\pi_\text{ref}(y \mid x)$.
\end{theorem}

\begin{sproof}
    Take any reward function $r(x, y)$, which defines an associated optimal model $\pi_r(y \mid x)$ as in Eq. \ref{eq:op_policy}. We will demonstrate that there exists a reward function within the equivalence class of $r$ that admits the stated reparameterization. Define the projection $f$ as  
\begin{equation}
    f(r; \pi_\text{ref}, \beta)(x, y) = r(x, y) - \beta\log\sum_{y}\pi_\text{ref}(y\mid x)\exp\left(\frac{1}{\beta}r(x, y)\right)
\end{equation}
This operator $f$ normalizes the reward function by subtracting the logarithm of the partition function associated with $\pi_r$. Because the normalization term depends only on the prefix $x$, $f(r; \pi_\text{ref}, \beta)(x, y)$ remains within the equivalence class of $r(x, y)$. Substituting $r$ on the right-hand side of Eq.~\ref{eq:main_eq} (which holds for any reward function) yields $f(r; \pi_\text{ref}, \beta)(x, y) = \beta \log \frac{\pi_r(y\mid x)}{\pi_\text{ref}(y\mid x)}$. Hence, $f$ maps $r$ to an equivalent reward function in the desired form, confirming that no generality is lost by adopting this reparameterization.
\end{sproof}

An alternative perspective on Theorem~\ref{thm:main} is that it precisely identifies which reward function within each equivalence class the DPO reparameterization chooses, specifically the reward function that satisfies:

\begin{equation}\label{eq:lag_p}
     \sum_{y}\underbrace{\pi_\text{ref}(y\mid x)\exp\left(\frac{1}{\beta}r(x, y)\right)}_{=\pi(y\mid x)\text{, by Thm.~\ref{thm:main} reparam.}} = 1,
\end{equation}

meaning that $\pi(y\mid x)$ forms a valid probability distribution (all probabilities are positive and sum to unity). Yet, according to Eq.~\ref{eq:op_policy}, Eq.~\ref{eq:lag_p} represents the partition function for the optimal policy generated by the reward function $r(x, y)$. The central insight behind the DPO method is that by imposing certain constraints on the otherwise under-determined Plackett-Luce (and particularly Bradley-Terry) family of preference models, we maintain the full set of representable reward functions while ensuring that the optimal policy described in Eq.~\ref{eq:op_policy} becomes analytically solvable for every prompt $x$.

\subsection{Instability of Actor-Critic Algorithms} Our framework can also be utilized to analyze the instabilities found in standard actor-critic algorithms employed in RLHF, such as PPO. We follow the RLHF workflow and concentrate on the RL fine-tuning phase described in Section \ref{section:prelims}. We establish connections to the control-as-inference framework \cite{levine2018reinforcement} applied to the constrained RL problem specified in \ref{eq:RL}. Assuming a parameterized policy model $\pi_{\theta}(y\mid x)$, the goal is to minimize the KL divergence $\mathbb{D}_{\text{KL}}[\pi_{\theta}(y|x) \mid \mid \pi^*(y\mid x)]$, where $\pi^*$ is the optimal policy from Eq.~\ref{eq:optimum_model} induced by the reward function $r_{\phi}(y, x)$. Through some algebraic manipulation, this objective can be expressed as:

\begin{equation}\label{eq:AC}
    \max_{\pi_{\theta}}\mathbb{E}_{\pi_{\theta}(y\mid x)}\bigg[\underbrace{r_{\phi}(x, y) -\beta\log\sum_{y}\pi_\text{ref}(y\mid x)\exp\left(\frac{1}{\beta}r_{\phi}(x, y)\right)}_{f(r_{\phi}, \pi_\text{ref}, \beta)} - \underbrace{\beta\log\frac{\pi_{\theta}(y\mid x)}{\pi_\text{ref}(y\mid x)}}_{\text{KL}}\bigg]
\end{equation}

This objective matches the one optimized in earlier studies \citep{ziegler2020finetuning, stiennon2022learning, bai2022training, ouyang2022training} by employing the DPO-equivalent reward for the reward class $r_{\phi}$. In this framework, the normalization component in $f(r_{\phi}, \pi_\text{ref}, \beta)$ can be seen as the soft value function of the reference policy $\pi_\text{ref}$. Although this term does not influence the optimal solution, omitting it may cause the policy gradient of the objective to exhibit high variance, thereby destabilizing learning. While it is possible to handle the normalization term using a learned value function, this approach can also present optimization challenges. Alternatively, prior research has normalized rewards by using a human completion baseline, which essentially acts as a single sample Monte-Carlo estimate of the normalization term. In contrast, the DPO reparameterization produces a reward function that eliminates the need for any baselines.

\section{Experiments}
In this section, we conduct empirical evaluations of DPO’s capability to train policies directly from preferences. Initially, in a controlled text-generation environment, we investigate how efficiently DPO balances maximizing reward against minimizing the KL-divergence with the reference policy, compared to standard preference learning methods such as PPO. Following this, we assess DPO’s performance on larger-scale models and more challenging RLHF tasks, including summarization and dialogue. We observe that with minimal hyperparameter tuning, DPO generally matches or surpasses strong baselines such as RLHF with PPO, as well as outperforms methods that select the best trajectory among $N$ samples under a learned reward function. Before detailing these findings, we outline the experimental setup; further information is provided in Appendix~\ref{app:exp_details}.

\textbf{Tasks.} Our study investigates three distinct open-ended text generation challenges. Across all experiments, the algorithms learn a policy from a preference dataset denoted as $\mathcal{D}=\bigl\{x^{(i)}, y_w^{(i)}, y_l^{(i)}\bigr\}_{i=1}^N$. In the case of \textbf{controlled sentiment generation}, $x$ represents the beginning segment of a movie review sourced from the IMDb dataset \cite{maas-EtAl:2011:ACL-HLT2011}, and the policy is tasked with producing $y$ that reflects a positive sentiment. To enable a controlled evaluation, this experiment \textit{generates} preference pairs by applying a pre-trained sentiment classifier such that $p(\text{positive}\mid x,y_w)>p(\text{positive}\mid x,y_l)$. For SFT, GPT-2-large is fine-tuned until convergence on the training split of the IMDb reviews (additional information is provided in App~\ref{app:sentiment_details}). In the \textbf{summarization} task, $x$ is a forum post from Reddit, and the policy must create a summary $y$ that captures the main ideas of the post. Consistent with prior research, we utilize the Reddit TL;DR summarization dataset \citep{volske-etal-2017-tl} along with human preference data collected by \citeauthor{stiennon2022learning}. The SFT model is fine-tuned on human-written summaries of forum posts\footnote{\url{https://huggingface.co/CarperAI/openai_summarize_tldr_sft}} using the TRLX \citep{leandro_von_werra_2023_7790115} framework for RLHF. The human preference dataset was compiled by \citeauthor{stiennon2022learning} from outputs generated by a different, but similarly fine-tuned, SFT model. Lastly, in the \textbf{single-turn dialogue} task, $x$ consists of a human query that can range from astrophysics questions to requests for relationship advice. The policy must generate a response $y$ that is both engaging and helpful. For this, we employ the Anthropic Helpful and Harmless dialogue dataset \citep{bai2022training}, which includes 170k conversations between humans and an automated assistant. Each conversation concludes with two response options produced by a large (though unspecified) language model, accompanied by a preference label indicating the human-favored response. Since there is no pre-trained SFT model available for this setting, we fine-tune a standard language model exclusively on the preferred completions to create the SFT model.

\textbf{Evaluation.} Our experiments employ two distinct evaluation strategies. To assess how well each algorithm optimizes the constrained reward maximization objective, in the controlled sentiment generation scenario we measure each algorithm based on the frontier of achieved reward and KL-divergence relative to the reference policy; this frontier is computable since the ground-truth reward function (a sentiment classifier) is accessible. Conversely, in practical applications where the ground truth reward function is unknown, we assess algorithms by their \textit{win rate} against a baseline policy, using GPT-4 as a surrogate for human evaluation of summary quality and response helpfulness in the summarization and single-turn dialogue settings, respectively. For summarization, baseline comparisons use reference summaries from the test set; for dialogue, the baseline is the preferred response found in the test data. Although prior research indicates that LMs can serve as more effective automated evaluators than standard metrics \citep{Chen2023ExploringTU}, we perform a human study to validate GPT-4’s use for evaluation in Sec.~\ref{sec:human-judgments}. We observe that GPT-4’s judgments correlate strongly with human assessments, with human agreement with GPT-4 generally comparable to or exceeding inter-annotator agreement.

\textbf{Methods.} Beyond DPO, we evaluate multiple existing methods for training language models to align with human preferences. At the simplest level, we test zero-shot prompting with \textbf{GPT-J} \citep{gpt-j} in summarization and 2-shot prompting with \textbf{Pythia-2.8B} \citep{biderman2023pythia} in dialogue. Additionally, we assess the \textbf{SFT} model and \textbf{Preferred-FT}, a model fine-tuned via supervised learning on the chosen completion $y_w$ from either the SFT model (in controlled sentiment and summarization) or a generic LM (in single-turn dialogue). Another pseudo-supervised technique is \textbf{Unlikelihood}~\citep{welleck2019neural}, which optimizes the policy to maximize the probability of $y_w$ while \textit{minimizing} the probability of $y_l$; an optional coefficient $\alpha\in[0,1]$ controls the weight on the unlikelihood term. We further examine \textbf{PPO} \citep{schulman2017proximal} trained with a reward function learned from preference data, and \textbf{PPO-GT}, an oracle using the ground truth reward function available in the controlled sentiment setup. For our sentiment experiments, we implement two versions of PPO-GT: an off-the-shelf version \cite{leandro_von_werra_2023_7790115} and a modified version that normalizes rewards and fine-tunes hyperparameters to enhance performance (these adjustments are also applied when running standard PPO with learned rewards). Lastly, we consider the \textbf{Best of $N$} baseline, which samples $N$ responses from the SFT model (or Preferred-FT in dialogue) and selects the highest-scoring response according to a reward model trained on preference data. This strong-performing approach separates the reward model’s quality from the PPO optimization process but is computationally expensive for even moderate $N$, as it requires generating $N$ completions for each query at test time.

\subsection{To what extent can DPO optimize the RLHF objective?}

The KL-constrained reward maximization objective common in RLHF algorithms seeks to balance maximizing reward with limiting how much the policy diverges from the reference policy. Hence, when evaluating different algorithms, it is important to consider both the reward obtained and the KL divergence; obtaining a marginally higher reward at the cost of significantly increased KL divergence may not be preferable. Figure~\ref{fig:frontier-tldr-main} depicts the reward-KL tradeoff frontier for several algorithms in the sentiment task. We conduct multiple training trials per algorithm, varying the policy conservativeness hyperparameters across runs (target KL $\in\{3,6,9,12\}$ for PPO, $\beta \in \{0.05,0.1,1,5\}$, $\alpha\in\{0.05,0.1,0.5,1\}$ for unlikelihood, and different random seeds for preferred-FT). This set comprises a total of 22 runs. After every 100 training steps until convergence, each policy is evaluated on a test prompt set, calculating the average reward under the true reward function alongside the average sequence-level KL\footnote{This refers to the sum of the per-timestep KL divergences.} with respect to the reference policy, $\text{KL}\left(\pi \mid \mid \pi_\text{ref}\right)$. Our findings indicate that DPO yields the most efficient frontier by a wide margin, attaining the highest reward while maintaining low KL divergence. This outcome is notable for several reasons. First, although DPO and PPO optimize the identical objective, DPO demonstrates significantly greater efficiency; its reward/KL tradeoff consistently surpasses that of PPO. Second, DPO surpasses PPO's frontier even when PPO is provided with the true reward signals (PPO-GT).

\subsection{Is DPO scalable to real-world preference datasets?}
\label{sec:dpo-real-datasets}
We proceed to assess the fine-tuning effectiveness of DPO on tasks involving summarization and single-turn dialogue. In summarization, automatic metrics like ROUGE often show weak alignment with human judgments~\citep{stiennon2022learning}, and previous studies have demonstrated that fine-tuning language models with PPO on human preference data yields more effective summaries. To evaluate various approaches, we generate completions on the test portion of the TL;DR summarization dataset and calculate the average win rate against reference completions from the test set. For all methods, completions are sampled at temperatures ranging from 0.0 to 1.0, with the resulting win rates illustrated in Figure~\ref{fig:frontier-tldr-main} (right). DPO, PPO, and Preferred-FT are all applied to fine-tune the same GPT-J SFT model\footnote{\url{https://huggingface.co/CarperAI/openai_summarize_tldr_sft}}. Our findings show that DPO attains a win rate around 61\% at temperature 0.0, surpassing PPO’s performance of roughly 57\% at its optimal temperature of 0.0. Additionally, DPO reaches a higher peak win rate compared to the best-of-$N$ baseline. We acknowledge that we did not extensively tune the $\beta$ hyperparameter in DPO, so these outcomes might underestimate its full capabilities. Furthermore, DPO demonstrates greater robustness to variations in sampling temperature compared to PPO, whose performance can deteriorate to match that of the base GPT-J model at elevated temperatures. Preferred-FT does not show notable improvements over the SFT model. A direct comparison between DPO and PPO via human evaluations is provided in Section~\ref{sec:human-judgments}, where DPO samples at temperature 0.25 were preferred 58\% of the time relative to PPO samples at temperature 0.

For the single-turn dialogue task, we assess various approaches using a subset of the test portion of the Anthropic HH dataset \citep{bai2022training}, involving one round of human-assistant interaction. GPT-4 evaluations utilize the preferred completions from the test set as references to calculate the win rates of different methods. Since there is no established SFT model for this task, we begin with a pre-trained Pythia-2.8B model, apply Preferred-FT to train a reference model on the selected completions ensuring that the completions align with the model’s distribution, and subsequently train using DPO. We also compare against the best result from 128 Preferred-FT completions (noting that the Best of $N$ baseline plateaus at 128 completions for this task; refer to Appendix Figure~\ref{fig:best-of-n}) and a 2-shot prompted variant of the Pythia-2.8B base model. Our findings indicate that DPO matches or exceeds performance at the optimal temperatures for each approach. Additionally, we evaluate an RLHF model trained with PPO on the Anthropic HH dataset \footnote{\url{https://huggingface.co/reciprocate/ppo_hh_pythia-6B}} from a reputable source \footnote{\url{https://github.com/CarperAI/trlx/tree/main/examples/hh}}, but we fail to identify a prompt or sampling temperature that surpasses the baseline Pythia-2.8B model’s performance. Drawing on our TL;DR results and given that both methods optimize the identical reward function, we treat the Best of 128 as an approximate benchmark for PPO-level performance. In summary, DPO stands as the sole computationally efficient method that enhances performance beyond the preferred completions in the Anthropic HH dataset, achieving comparable or superior outcomes relative to the computationally intensive Best of 128 baseline. Lastly, Figure~\ref{fig:dialogue-main} demonstrates that DPO attains its peak performance relatively rapidly.

\subsection{Generalization to a new input distribution}

To further assess the robustness of PPO and DPO under shifts in data distribution, we test the PPO and DPO policies from our Reddit TL;DR summarization experiment on a distinct distribution: news articles from the test portion of the CNN/DailyMail dataset \citep{nallapati-etal-2016-abstractive}, employing the optimal sampling temperatures found for TL;DR (0 and 0.25). The outcomes are summarized in Table~\ref{tab:ood}. We measured the GPT-4 win rate against the ground-truth summaries in these datasets, using the identical GPT-4 (C) prompt applied in Reddit TL;DR, but substituting the phrase ``forum post'' with ``news article''. On this new distribution, DPO maintains a substantial performance advantage over PPO. This experiment offers preliminary evidence that DPO policies can generalize on par with PPO policies, despite DPO not utilizing the extra unlabeled Reddit TL;DR prompts that PPO leverages.

\subsection{Validating GPT-4 judgments with human judgments}
\label{sec:human-judgments}
We carry out a human evaluation study to confirm the trustworthiness of GPT-4’s judgments, utilizing data from the TL;DR summarization experiment with two distinct GPT-4 prompts. The \textbf{GPT-4 (S)} (simple) prompt straightforwardly asks which summary better captures the important information in the post. The \textbf{GPT-4 (C)} (concise) prompt additionally requests which summary is more concise; this prompt was tested because we observed that with the \textbf{GPT-4 (S)} prompt, GPT-4 tends to favor longer, more repetitive summaries than humans typically do. Full prompts can be found in Appendix~\ref{app:prompts}. We conduct three pairwise comparisons involving the highest-performing method (DPO, temp. 0.25), the lowest-performing (PPO, temp. 1.0), and an intermediate performer (SFT, temp. 0.25), aiming to encompass a range of sample quality; each method is compared against greedily-sampled PPO (its best-performing temperature). Results indicate that GPT-4 agrees with human judgments about as frequently as humans agree among themselves under both prompts, implying that GPT-4 serves as a reasonable stand-in for human evaluations (due to a limited number of human raters, multiple human judgments were collected only for the DPO and PPO-1 comparisons). Overall, the \textbf{GPT-4 (C)} prompt generally yields win rates that more closely align with human judgments; consequently, we employ this prompt for the main analyses in Section~\ref{sec:dpo-real-datasets}. Additional information about the human study, including the rater interface and the list of human participants, is available in Appendix~\ref{app:human-study}.

\section{Discussion}
Learning from preferences offers a robust and scalable approach for developing capable and aligned language models. We have presented DPO, a straightforward training method that enables language models to be trained from preferences without resorting to reinforcement learning. Instead of forcing the preference learning task into a conventional RL framework to leverage existing RL algorithms, DPO establishes a direct correspondence between language model policies and reward functions, allowing language models to be trained to meet human preferences \textit{directly} using a simple cross-entropy loss, without involving reinforcement learning or compromising generality. With minimal hyperparameter tuning, DPO achieves performance comparable to or exceeding that of current RLHF techniques, including those based on PPO; consequently, DPO significantly lowers the barrier to training language models using human preference data.

\textbf{Limitations \& Future Work.} Our findings prompt several key questions for further investigation. How well does the DPO policy generalize beyond the training distribution compared to models trained with explicit reward functions? Preliminary evidence indicates that DPO policies generalize on par with PPO-based approaches, yet a more thorough analysis is necessary. For instance, can DPO-based training with self-labeling effectively leverage unlabeled prompts? Additionally, how does reward over-optimization present itself within the direct preference optimization framework, and might the slight performance dip shown in Figure~\ref{fig:dialogue-main}-right be an example of this phenomenon? Furthermore, although our evaluations cover models up to 6B parameters, scaling DPO to much larger, state-of-the-art models is a promising avenue for future research. Regarding evaluations, we observe that GPT-4’s win rate assessments are influenced by the prompt; subsequent work could explore optimal strategies for eliciting high-quality judgments from automated evaluators. Lastly, numerous potential applications for DPO extend beyond preference-based language model training, including training generative models across various modalities.

\end{document}